% Source: https://artofproblemsolving.com/community/c3h1668907_cmms_problem_3340_unclarity
% Stable identity: aops-topic-1668907-q01
A product is odd exactly when all three digits are odd. There are only five odd digits, $1,3,5,7,9$, so at most one set can consist of odd digits. Hence ``at least one set'' and ``exactly one set'' with an odd product are the same event, and the answer is the same for both readings.

Label the sets first, second and third (this does not change any probability). The number of divisions is $\dbinom93\dbinom63=84\cdot20=1680$. For a favourable division, choose which set is all odd ($3$ ways) and its digits ($\dbinom53=10$ ways), then split the remaining six digits into the other two sets ($\dbinom63=20$ ways): $3\cdot10\cdot20=600$ divisions. The probability is
\[
\dfrac{600}{1680}=\boxed{\dfrac{5}{14}},
\]
and it does not change for ``exactly one''.
